\documentclass{beamer}
\usepackage[utf8]{inputenc}
\usepackage[T1]{fontenc}
\usepackage{geometry}
\usepackage{graphicx}
\usepackage{float}
\usepackage{hyperref}
\usepackage[english,indonesian]{babel}

\title{Optimasi Pemrosesan Template LaTeX Menggunakan Pendekatan Berbasis AI}
\author{Academic Support Agent}
\date{\today}

\begin{document}
\maketitle

\begin{abstract}

\end{abstract}

\chapter{PENDAHULUAN}


\chapter{TINJAUAN PUSTAKA}


% Original content (commented out):
% \documentclass{beamer}

\usetheme[compactlogo,white]{Madison}
\usefonttheme[onlymath]{serif}


\usepackage{amsmath}
\usepackage{amssymb,amsfonts,amsthm}
\usepackage{graphicx}
\graphicspath{{figs/}}
\usepackage{natbib}


\title{A Cool Title}
\subtitle{a clever subtitle}

\author{First A.\ Author and Second B. Author}

\date[]{
  \\
  Workshop on
  \emph{Various Mathematical Things}\\
  Cool Institute, Somewhere  (Online)\\
  1 January 2098
}

\begin{document}

\small

\begin{frame}[plain]
  \maketitle
\end{frame}



\begin{frame}{A slide}

  \begin{center}
    \includegraphics[width=.5\textwidth]{decay_corr_sketch}
  \end{center}

  \begin{enumerate}[<+->]
  \item Item 1
  \item Item 2 and a citation \cite{Thurston1988}.
\end{enumerate}
\end{frame}



\bibliographystyle{plainnat}

\begin{frame}[t,allowframebreaks]
  \frametitle{References}
  \bibliography{main}
\end{frame}

\end{document}


\end{document}